\documentclass[aspectratio=169,11pt]{beamer}
\usepackage[T1]{fontenc}
\usepackage{lmodern}
\usepackage{graphicx}
\usepackage{booktabs}
\usepackage{array}
\newcommand{\DPTitle}{Dual-threshold time-walk calibration for CMS BTL prototypes}
\newcommand{\DPShortTitle}{Dual-threshold time-walk calibration}
\newcommand{\DPContact}{cms-dpg-conveners-mtd@cern.ch}
\newcommand{\DPAbstract}{%
A charge-dependent time-walk model is studied for CMS barrel timing layer
(BTL) prototype channels using two discriminator thresholds. Their current
separation divided by the measured crossing-time separation provides a
finite-interval pulse-slope proxy. A linear profile fit describes this
proxy as a function of the measured charge, and an inverse-slope function
models the timing correction with one amplitude parameter per channel.

Slope measurements from run 4076 are compared with independent
microchannel-plate (MCP) reference calibrations from run 4399. The two runs
contain different events. A free constant offset and the correction amplitude
are fitted only where both measurements have populated charge support.
The resulting agreement is a calibration-shape comparison, not an independent
measurement of timing resolution.

The draft also examines channel-to-channel variations and an exploratory
prediction of the correction amplitude from internal pulse information.
Such a prediction could enable operation without an event-by-event external
time reference, but its transferability and full uncertainty remain to be
established.}

\newcommand{\AnalysisRepo}{https://github.com/De-Cristo/LecroyTranslator4CMSTimingDetector}
\newcommand{\AnalysisRevision}{45a907c422d6c8afbe0cdf44cd2827d3433d0f97}
\newcommand{\RepoFile}[2]{\href{\AnalysisRepo/blob/\AnalysisRevision/#1}{#2}}
\usetheme{default}
\setbeamercolor{normal text}{fg=black,bg=white}
\setbeamercolor{structure}{fg=black}
\setbeamercolor{frametitle}{fg=black,bg=white}
\setbeamerfont{frametitle}{size=\large,series=\bfseries}
\setbeamersize{text margin left=7mm,text margin right=7mm}
\setbeamertemplate{navigation symbols}{}
\setbeamertemplate{headline}{%
  \hbox to \paperwidth{\hspace{7mm}\scriptsize\color{black!55}Internal review draft -- 22 October 2026
  \hfill\raisebox{0pt}[9mm][1mm]{\includegraphics[height=7mm]{CMSlogo_black_label_1024.png}}\hspace{6mm}}}
\setbeamertemplate{footline}{%
  \hbox to \paperwidth{\hspace{7mm}\scriptsize\DPShortTitle
  \hfill Draft v1\quad\insertframenumber/\inserttotalframenumber\hspace{7mm}}\vspace{2mm}}
\setbeamertemplate{frametitle}{\vspace{-1mm}\insertframetitle\par\vspace{1mm}}
\setlength{\parskip}{2.5pt}
\newcommand{\Review}[2]{\par\smallskip{\fontsize{7.5}{8.5}\selectfont\color{red!65!black}\textbf{Review #1:} #2\par}}
\newcommand{\Caption}{\fontsize{9}{10.5}\selectfont}
\newcommand{\FigureSlide}[4]{%
  \begin{frame}[t]{#1}
    \begin{columns}[T,onlytextwidth]
      \begin{column}{.61\textwidth}
        \centering\includegraphics[width=\linewidth,height=52mm,keepaspectratio]{figures/#2.pdf}
      \end{column}
      \begin{column}{.37\textwidth}\Caption #3\par\medskip{\fontsize{7}{8.5}\selectfont\color{black!55}#4\par}\end{column}
    \end{columns}
  \end{frame}}
\newcommand{\OpeningSlides}{%
  \begin{frame}
    \centering{\LARGE\DPTitle\par}\vspace{7mm}
    {\large CMS\par}\vspace{3mm}
    \href{mailto:\DPContact}{\DPContact}
    \vspace{7mm}\par{\small Discussion draft; not approved for public release.}
  \end{frame}
  \begin{frame}{Abstract}\small\DPAbstract\end{frame}}
\hypersetup{pdftitle={\DPTitle},pdfauthor={CMS}}

% Generated from frozen numerical results by build_assets.py.
\newcommand{\TimingCount}{10127}
\newcommand{\DBefore}{55.16}
\newcommand{\DAfter}{54.82}
\newcommand{\DLow}{54.32}
\newcommand{\DHigh}{55.22}
\newcommand{\DSD}{61.26}
\newcommand{\DTail}{1.57}
\newcommand{\RBefore}{77.69}
\newcommand{\RAfter}{64.16}
\newcommand{\RLow}{63.27}
\newcommand{\RHigh}{64.99}
\newcommand{\RSD}{70.33}
\newcommand{\RTail}{1.56}
\newcommand{\RefBinsA}{19}
\newcommand{\RefBinsB}{19}
\newcommand{\RefBinsC}{20}
\newcommand{\RefBinsD}{19}
\newcommand{\ExampleAWC}{18.332}
\newcommand{\ExampleLow}{18.120}
\newcommand{\ExampleHigh}{18.607}
\newcommand{\ExampleRMS}{6.80}
\newcommand{\GapCount}{703}
\newcommand{\SummaryChi}{538.6}
\newcommand{\PredictorRMSE}{3.22}
\newcommand{\PredictorFraction}{14.1}

\begin{document}
\OpeningSlides

\begin{frame}{Two complementary runs, not the same events}
\small
\begin{tabular}{l>{\raggedright\arraybackslash}p{4.8cm}>{\raggedright\arraybackslash}p{5.2cm}}
\toprule
Sample & Available information & Role in this study\\
\midrule
Run 4076 & Valid dual-threshold timing and charge & Measure the pulse-slope proxy $S(Q)$\\
Run 4399 & MCP-associated reference timing and charge & Measure the empirical walk $W_{\rm MCP}(Q)$\\
\bottomrule
\end{tabular}\par
\medskip
The runs contain different events; no event-by-event matching is performed.
The comparison transfers a channel-response model between runs.

The upstream and downstream modules each contain bars read out at two ends.
Detector/readout context is given in Refs.~\cite{MTDTDR,TOFHIR2}.
The channel-to-bar mapping determines the coincidence pairs; the same channel
number is compared across the two runs.

The comparison assumes the same nominal detector setup in both runs,
with MCP reference timing in run 4399 and dual-threshold timing in run 4076.
Cross-run response differences are not included in the quoted uncertainties.
\end{frame}

\begin{frame}{From two crossing times to a walk model}
\small
For thresholds $I_1<I_2$ and their valid crossing times $t_1,t_2$, define
\[
 S=\frac{I_2-I_1}{t_2-t_1}\quad [\mu{\rm A/ns}],\qquad
 \overline S(Q)\simeq c_0+c_1Q.
\]
$S$ is a finite-interval rising-edge slope proxy. It is not the current
difference $\Delta I$ itself and is not necessarily the instantaneous
derivative at the first threshold.
\[
 W_Q(Q)=\frac{A}{c_0+c_1Q},\qquad t_{\rm corr}=t_{\rm measured}-W_Q(Q).
\]
The amplitude $A$ (AWC) has units of current and is constant for a channel
and operating configuration. The correction varies event by event through $Q$.
\end{frame}

\begin{frame}{Timing reconstruction and threshold scale}
\small
The run-4076 sample is file \texttt{23\_e.root}. For coarse codes $C_1,C_2$,
\[
 d_c=(C_2-C_1)\bmod1024,\qquad
 t_2=t_1+(6250\,\mathrm{ps})d_c-f_1+f_2.
\]
Subtraction uses signed 64-bit coarse codes. The coarse-time term uses
floating-point arithmetic to avoid the previous 16-bit multiplication overflow.

The fine-time lookup stores the mean of $t_1\bmod6250\,\mathrm{ps}$ for each
(channel, TAC, fine-code) combination, using all input hits. This lookup supplies
$f_1$ and $f_2$, with entries for every selected hit in the 22 audited channels.

All selected hits have threshold register \texttt{step2=113621}, decoded as
DAC codes 10 and 35. The implemented scale gives
\[
 \Delta I=(35-10)\times0.3125\,\mu\mathrm{A}=7.8125\,\mu\mathrm{A},\qquad
 \frac{S}{\mu\mathrm{A/ns}}=1000\,\frac{\Delta I/\mu\mathrm{A}}{(t_2-t_1)/\mathrm{ps}}.
\]
{\scriptsize \RepoFile{Q_Walk/qwalk_core.py}{GitHub: timing and threshold reconstruction}.
This verifies the implemented conversion, not an independent calibration of the hardware thresholds.}
\end{frame}

\begin{frame}{Selection and profile construction}
\small
\begin{itemize}
\item Require one hit at each end of the target bar and no other hit in that
module. The other module is unconstrained; channel 192 is ignored.
\item Require finite $Q,S$, $\Delta I>0$, $t_2-t_1>0$ and $d_c=0$.
This selects the same-coarse-bin prompt population.
\item Require $Q\geq Q_{\rm left}$, the left half-height boundary of the
smoothed charge histogram. The Landau--Gaussian fit does not set this cut.
There is no additional upper-charge cut or 100\,ps minimum-rise cut.
\item Form ten equal-width bins over the selected charge range. Plot mean $S$
and hit SD; fit with SEM weights. Bins with undefined or negligible SEM are omitted.
\item Bootstrap selected hits with fixed bin edges and refit 1000 times.
Quote the 16th--84th percentiles.
\end{itemize}
{\scriptsize All selected hits are used; there is no training or test split.
\RepoFile{Q_Walk/qwalk_two_hit_pair_profile.py}{GitHub: pair selection and profiling},
with the newer optional minimum-rise cut disabled.}
\end{frame}

\begin{frame}{Reproducing the frozen run-4076 profiles}
\small
Reprocessing \texttt{23\_e.root} reproduces the selected counts, ten-bin profiles
and linear-fit coefficients for all 22 channels. The saved bootstrap coefficients
reproduce the reported percentile intervals.
\begin{center}
\begin{tabular}{lrrr}
\toprule
Channel & Target-pair hits & Prompt hits & Peak-right hits\\
\midrule
137 & 25979 & 20789 & 19126\\
\bottomrule
\end{tabular}
\end{center}
For channel 137, the histogram-derived boundary is $Q_{\rm left}=184.9163$ ADC.
Each profile point is placed at the mean measured $Q$ in its bin,
not at the geometric bin center.

The frozen selection retains rises below 100\,ps in channels 130, 139, 151,
159, 236 and 238. These events remain in the sample. Their effect
requires selection-sensitivity and channel-quality studies.

Channel 135 has one bin with a single entry. Its SEM is undefined, so only nine
bins enter the linear fit. The audit changed neither the fits nor the plotted data.

\end{frame}

\begin{frame}{Reference version shared with the MCP reference note~\cite{ReferenceNote}}
\small
The run-4399 MCP functions were refitted on 22 October 2026. The updated procedure
detects segment boundaries across missing trigger timestamps and excludes invalid
clock groups. The MCP reference note documents the timing validation and remaining limitations.

For each topology, file and segment, the scaled MAD of the channel-192
versus scope-trigger line residuals must not exceed 1\,ns before LOWESS~\cite{LOWESS}.
The selection uses clock information, not SiPM timing residuals.
Suspected incorrect event associations are excluded; MCP rows are not shifted.

All AWC-and-offset fits in this note use these updated references.
The run-4076 slope fits, bootstrap replicas and measured-support matching
definition are unchanged. The charge spectra use the updated reference sample.

The plotted AWC intervals hold the reference fixed. They exclude reference-fit
uncertainty, clock-selection uncertainty and uncertainty in transfer between runs.
The uncertainty table later in this note adds reference-fit statistics separately.

\end{frame}

\FigureSlide{Charge spectra define the available comparison sample}{q_spectra_ch137}{
The spectra show the channel-137 charge distributions in the two independent runs.
Lines show hits passing the quality selection; shaded spectra also require
the charge-peak selection. Each run is normalized to its own quality-selected count.

The charge proxy is in ADC units, not MeV. The samples differ in population and
acceptance, so the curve alignment uses only charge intervals populated in both runs.
}{Recovered histogram counts, not a new peak fit.
Run-4076 profile selection starts at about 184.9 ADC for this channel.}

\FigureSlide{Channel 133: fitting the slope profile}{slope_ch133}{
Run 4076 channel-133 hits are grouped into ten charge bins.
Points show mean $S$; bars show the SD of individual selected hits.
The line is the SEM-weighted fit $c_0+c_1Q$.

The shaded interval is the central 68\% bootstrap interval of the fitted
mean curve, shown without enlargement.
}{SD is the hit standard deviation; SEM is the standard error of the mean.
The band describes uncertainty on the fitted mean curve.}

\FigureSlide{Channel 154: the opposite end of bar 7}{slope_ch154}{
Channel 154 uses the same selection and profiling procedure.
The mean slope increases with charge, but the response is calibrated
independently at each bar end.

Bars show the hit SD. The line uses SEM weights. The band propagates the
bootstrap coefficient replicas, including their correlation.
Linearity is a working approximation, not a validated precision model.
Model and selection dependence are excluded from the quoted errors;
the separate sensitivity study is retained outside this note.
}{Run 4076, saved profile result, displayed without refitting.}

\FigureSlide{Channel 137: pulse slope and charge}{slope_ch137}{
The selected run-4076 sample contains 19\,126 channel-137 hits.
Ten profile points constrain a linear pulse-slope model over the observed
charge range.

The SD bars show event-to-event spread, while SEM weights determine the fit.
The narrow shaded band is the conditional 68\% bootstrap interval of the
mean prediction, not a single-hit uncertainty.
}{The count and fit come from the saved profile results.
Uncertainty bands are shown without enlargement.}

\FigureSlide{Channel 150: the opposite end of bar 8}{slope_ch150}{
Channel 150 uses the same profiling procedure.
The charge-dependent slope provides the denominator of the proposed walk
function; a separate amplitude $A$ determines the correction strength.

The line and shaded interval describe the fitted mean. Points and SD bars
show the spread within the selected charge bins.
}{Run 4076, saved profile result. A constant AWC gives a
charge-dependent timing correction.}

\begin{frame}{Determine AWC using only jointly measured support}
\small
The run-4399 reference is a fitted quadratic in inverse charge,
\[
 W_{\rm MCP}(Q)=a+b/Q+c/Q^2.
\]
Let $\Omega$ be the intersection of the two outer profile-mean ranges,
restricted to charge bins populated in both runs. Fit
\[
 (A,B)=\arg\min_{A,B}\sum_{Q_k\in\Omega}
 \left[\frac{A}{c_0+c_1Q_k}+B-W_{\rm MCP}(Q_k)\right]^2.
\]
The fit samples a uniform 3000-point charge grid, then removes intervals
without data support. The offset $B$ absorbs the arbitrary time zero.

This is uniform-charge shape matching, not an event-weighted likelihood.
It uses neither extrapolated tails nor a single central anchor point.
The uncertainty study below evaluates sensitivity to charge weighting and
correlated MCP-fit uncertainty. Uniform Q is the nominal choice;
it has not been established as uniquely optimal.
\end{frame}

\FigureSlide{Channel 137: inverse-slope and MCP walk shapes}{comparison_ch137}{
The black curve is the run-4399 MCP quadratic. The red curve is the
run-4076 inverse-slope model after fitting AWC and a constant offset.
The curves are shown only in charge intervals populated in both runs.

Lower panel: run-4076 model minus MCP model. Its RMS shape mismatch is
\ExampleRMS\,ps on the supported uniform-Q grid, not an event timing resolution.

The band varies the run-4076 slope fit while holding the MCP curve fixed.
The uncertainty table later in this note also includes MCP-fit uncertainty.
}{AWC = \ExampleAWC\,$\mu$A; conditional 16--84\% interval
[\ExampleLow, \ExampleHigh]\,$\mu$A. Both curves share one display offset.}

\FigureSlide{Channel 154: a second shape-comparison example}{comparison_ch154}{
The same objective, restricted to the measured overlap, determines AWC and the offset for
channel 154. Both plotted curves are calibrated models, not matched
events from different runs.

The lower panel shows the remaining shape mismatch, with
no additional line subtracted. This measures model-shape closure,
not corrected-event timing resolution.
}{Conditional bootstrap band: 4076 slope replicas, refitted AWC and
offset, fixed 4399 quadratic.}

\FigureSlide{Channel 159: missing support is not extrapolation}{comparison_ch159}{
Only intervals with populated profile bins in both runs contribute to
the alignment. Breaks in the displayed curves mark excluded charge intervals.
They do not represent discontinuities in either analytic model.

The lower panel uses the same supported points. All displayed
intervals participate in the common AWC-and-offset fit.
It is a model difference, not an event-time distribution.
}{The comparison excludes \GapCount{} of the 3000 initial grid points
for this channel. No line is drawn through those gaps.}

\FigureSlide{Downstream example: channel 242}{comparison_ch242}{
This example applies the method to a downstream channel.
In run 4399, upstream bar 8 supplies the coincidence pair alongside
the downstream target pair and MCP indicator.

In run 4076, the target downstream pair is required, without fixing
an upstream pair. The curves use only their common measured support.
Downstream comparisons also assume the same nominal detector setup in both runs.
Selection-dependent response differences remain unquantified.
}{The lower panel is a model difference, not a timing-resolution measurement.
The other module is not used as an internal timing reference to determine AWC here.}

\FigureSlide{A selected subset does not establish a universal law}{awc_summary_selected}{
The ten-channel upstream subset previously selected for discussion is
shown on linear axes. Bars are slope-only statistical AWC bootstrap intervals;
channel numbers identify the points.

The weighted straight-line fit has $\chi^2/{\rm ndf}=\SummaryChi/8$ under its
limited error model. It does not adequately describe the variation.
The subset was selected interactively and is illustrative, not quality-selected
or an unbiased population-level validation.
}{The fit treats x as fixed and uses half the 16--84\% AWC span as
its symmetric uncertainty. Correlations and reference errors are omitted.}

\begin{frame}{Channel coverage and limitations}
\footnotesize
The explored set contains 14 upstream and 8 downstream channels, with 21
completed alignments. Only channel 238 lacks measured overlap.
No additional quality cut is applied to this summary.
\begin{center}
\begin{tabular}{l>{\raggedright\arraybackslash}p{9.4cm}}
\toprule
Channels & Observed limitation and treatment\\
\midrule
238 & No common measured charge range; no AWC estimate.\\
235, 244 & Overlap spans only 3.10 and 16.99 ADC, respectively;
one profile mean from each run lies within each overlap. Fits retained.\\
130, 236, 239 & Slope-only AWC intervals include zero and span hundreds
of $\mu$A. Fits retained as poorly constrained cases.\\
136, 245 & Negative central AWC values are retained, not accepted
as validated positive-amplitude calibrations.\\
\bottomrule
\end{tabular}
\end{center}
The ten-channel illustration is unchanged: 133, 134, 135, 137, 150, 151,
154, 155, 156 and 159. This subset does not define a channel-selection criterion.

{\scriptsize Table entries are diagnostic examples, not additional exclusions.
Small curve mismatch over narrow support does not validate a calibration.
The full per-channel support and interval table is saved with the figure data.}
\end{frame}

\begin{frame}{Uncertainties}
\small
\begin{tabular}{l>{\raggedright\arraybackslash}p{9.1cm}}
\toprule
Displayed quantity & Meaning\\
\midrule
Profile bars & SD of selected hits in each charge bin; not the uncertainty
on the mean.\\
Fit weights & SEM of each profile mean; not the SD.\\
Slope band & 16--84\% range from 1000 bootstrap linear fits with fixed bin edges.\\
AWC bars / walk band & Propagate the slope replicas through a new AWC-plus-offset
alignment to the same fixed reference curve.\\
\bottomrule
\end{tabular}\par
\medskip
These definitions refer to the original displayed bands and bars.
The next table additionally propagates the run-4399 quadratic covariance.
Both prescriptions exclude uncertainty in response transfer between runs,
clock and fine-time calibration, selection and model choices, and cross-channel correlations.
Neither is a total calibration uncertainty or a timing-resolution estimate.
\end{frame}

\begin{frame}{MCP-fit uncertainty and charge weighting}
\Caption
For each slope-bootstrap replica, propagate the correlated MCP coefficient
covariance through the AWC-plus-offset fit. The Gaussian mixture gives the
16--84\% interval; central AWC values and comparison curves are unchanged.
\par
All table values are in $\mu$A. ``Both runs'' adds MCP-fit statistics, not
cross-run or model systematics. The original plots retain slope-only statistical bands.
{\centering\begin{tabular}{rrrrr}
\toprule
Channel & AWC & Slope only & Both runs & Max. weight shift\\
\midrule
133 & 36.03 & [32.73, 44.75] & [32.48, 45.72] & 4.11\\
154 & 24.25 & [23.92, 24.86] & [22.98, 25.84] & 0.85\\
137 & 18.33 & [18.12, 18.61] & [17.91, 18.80] & 2.05\\
150 & 20.89 & [20.65, 21.15] & [20.41, 21.38] & 1.13\\
\bottomrule
\end{tabular}
\par}
Maximum weighting shift: largest $|\Delta\mathrm{AWC}|$ using uniform $1/Q$
or either run's charge population instead of uniform Q. The overlap and free
offset are unchanged. These shifts are not statistical errors and are not added in quadrature.
\par
All 21 overlap channels were checked; channel 238 has no overlap.
A shared reference-time offset affects only the fitted offset, not AWC.
These are conditional statistical intervals, with weighting sensitivity
reported separately, not total calibration uncertainties.
\end{frame}

\begin{frame}{Result and next validation step}
\small
\begin{itemize}
\item A dual-threshold measurement provides a charge-dependent slope proxy;
an inverse-slope function can be aligned to independent MCP walk curves.
\item Restricting alignment to jointly measured charge support avoids
determining AWC from extrapolation.
\item Reference-calibrated AWC varies between channels. A straight line
against $1/c_1$ is insufficient for the explored subset.
\item The present calibration still uses the MCP to determine AWC.
It does not establish reference-free AWC extraction by itself.
\item Before use without MCP, an internal predictor must be tested on held-out
bars and configurations, including uncertainties in transfer and in the target AWC values.
\end{itemize}
\medskip
The following backup explores such a predictor; it is not a final
performance claim or a demonstration with cosmic rays.
\end{frame}

\begin{frame}{Backup: an internal predictor for channel AWC}
\small
For $W=A/S$ and $S=c_0+c_1Q$,
\[
 -\frac{dW}{d\ln Q}=\frac{A\,c_1Q}{S(Q)^2}.
\]
If this local walk derivative is approximately shared between channels at
a prescribed typical charge $Q_*$, a candidate model is
\[
 A_{\rm pred}=K\,\frac{S(Q_*)^2}{c_1Q_*}.
\]
Here $Q_*$ is the median charge of the run-4076 quality- and peak-right-selected
sample. It is computed once per channel; $A_{\rm pred}$ remains constant
when correcting individual events.

The common $K$ is learned from MCP-calibrated training channels. Predictions
for new channels use internal measurements, but model development uses the reference.
The shared derivative and choice of $Q_*$ remain hypotheses requiring
independent validation across operating conditions, without tuning to the test results.
\end{frame}

\FigureSlide{Backup: predict a bar excluded from the model fit}{internal_predictor}{
Each point compares reference-fitted AWC with a prediction made after
excluding that channel's entire bar from the fit of $K$.
Ten upstream channels span six held-out bar groups.

The RMS prediction error is \PredictorRMSE\,$\mu$A; the RMS fractional error is \PredictorFraction\%.
Horizontal intervals are conditional reference-AWC intervals, not
predictive intervals.
This small selected sample leaves input--target correlations and full
prediction uncertainties unpropagated. It does not demonstrate reference-free calibration.
}{Exploratory median-$Q_*$ candidate. The dashed line indicates equality,
not a fitted result. A general 10\% precision claim is not established.}

\begin{frame}{References}
\fontsize{9}{10.5}\selectfont
\setbeamertemplate{bibliography item}[text]
\begin{thebibliography}{4}
\bibitem{MTDTDR}
CMS Collaboration, \emph{A MIP Timing Detector for the CMS Phase-2 Upgrade},
CERN-LHCC-2019-003, CMS-TDR-020 (2019).
\href{https://cds.cern.ch/record/2667167}{CERN document record}.
\bibitem{TOFHIR2}
E. Albuquerque et al., \emph{TOFHIR2: The readout ASIC of the CMS Barrel
MIP Timing Detector} (2024).
\href{https://arxiv.org/abs/2404.01208}{arXiv:2404.01208}.
\bibitem{LOWESS}
W. S. Cleveland, \emph{Robust Locally Weighted Regression and Smoothing
Scatterplots}, J. Am. Stat. Assoc. \textbf{74} (1979) 829--836.
\href{https://doi.org/10.1080/01621459.1979.10481038}{doi:10.1080/01621459.1979.10481038}.
\bibitem{Bootstrap}
B. Efron, \emph{Bootstrap Methods: Another Look at the Jackknife},
Ann. Statist. \textbf{7} (1979) 1--26.
\href{https://doi.org/10.1214/aos/1176344552}{doi:10.1214/aos/1176344552}.
\end{thebibliography}
\medskip
These sources provide detector/readout context and methodological background.
The run-specific selections, implementations and numerical results are those
of this study, not measurements or prescriptions taken from these references.
\end{frame}


\end{document}
